\chapter{2022年新高考I卷}

\section{单选题}

\begin{problem}
若集合 \(M=\{x\mid \sqrt{x}<4\}\)，\(N=\{x\mid 3x\ge 1\}\)，则 \(M\cap N=\)（\quad）

\choices
    {\(\{x\mid 0\le x<2\}\)}
    {\(\left\{x\mid \frac13\le x<2\right\}\)}
    {\(\{x\mid 3\le x<16\}\)}
    {\(\left\{x\mid \frac13\le x<16\right\}\)}
\end{problem}

\begin{problem}
若 \(\i(1-z)=1\)，则 \(z+\overline{z}=\)（\quad）

\choices
    {\(-2\)}
    {\(-1\)}
    {1}
    {2}
\end{problem}

\begin{problem}
在 \(\triangle ABC\) 中，点 \(D\) 在边 \(AB\) 上，\(BD=2DA\). 记 \(\overrightarrow{CA}=\bs m\)，\(\overrightarrow{CD}=\bs n\)，则 \(\overrightarrow{CB}=\)（\quad）

\choices
    {\(3\bs m-2\bs n\)}
    {\(-2\bs m+3\bs n\)}
    {\(3\bs m+2\bs n\)}
    {\(2\bs m+3\bs n\)}
\end{problem}

\begin{problem}
南水北调工程缓解了北方一些地区水资源短缺问题，其中一部分水蓄入某水库. 已知该水库水位为海拔 \(148.5\text{m}\) 时，相应水面的面积为 \(140.0\text{km}^2\)；水位为海拔 \(157.5\text{m}\) 时，相应水面的面积为 \(180.0\text{km}^2\). 将该水库在这两个水位间的形状看作一个棱台，则该水库水位从海拔 \(148.5\text{m}\) 上升到 \(157.5\text{m}\) 时，增加的水量约为（\(\sqrt7\approx 2.65\)）（\quad）

\choices
    {\(1.0\times 10^9\text{m}^3\)}
    {\(1.2\times 10^9\text{m}^3\)}
    {\(1.4\times 10^9\text{m}^3\)}
    {\(1.6\times 10^9\text{m}^3\)}
\end{problem}

\begin{problem}
从 \(2\) 至 \(8\) 的 \(7\) 个整数中随机取 \(2\) 个不同的数，则这 \(2\) 个数互质的概率为（\quad）

\choices
    {\(\dfrac16\)}
    {\(\dfrac13\)}
    {\(\dfrac12\)}
    {\(\dfrac23\)}
\end{problem}

\begin{problem}
记函数 \(f(x)=\sin\left(\omega x+\frac{\pi}{4}\right)+b\)（\(\omega>0\)） 的最小正周期为 \(T\). 若 \(\frac{2\pi}{3}<T<\pi\)，且 \(y=f(x)\) 的图象关于点 \(\left(\frac{3\pi}{2},2\right)\) 中心对称，则 \(f\left(\frac{\pi}{2}\right)=\)（\quad）

\choices
    {1}
    {\(\dfrac32\)}
    {\(\dfrac52\)}
    {3}
\end{problem}

\begin{problem}
设 \(a=0.1e^{0.1}\)，\(b=\frac{0.1}{1-0.1}\)，\(c=-\ln 0.9\)，则（\quad）

\choices
    {\(a<b<c\)}
    {\(c<b<a\)}
    {\(c<a<b\)}
    {\(a<c<b\)}
\end{problem}

\begin{problem}
已知正四棱锥的侧棱长为 \(l\)，其各顶点都在同一球面上. 若该球的体积为 \(36\pi\)，且 \(3\le l\le 3\sqrt3\)，则该正四棱锥体积的取值范围是（\quad）

\choices
    {\(\left[18,\frac{81}{4}\right]\)}
    {\(\left[\frac{27}{4},\frac{81}{4}\right]\)}
    {\(\left[\frac{27}{4},\frac{64}{3}\right]\)}
    {\([18,27]\)}
\end{problem}

\section{多选题}

\begin{problem}
已知正方体 \(ABCD-A_1B_1C_1D_1\)，则（\quad）

\choices
    {直线 \(BC_1\) 与 \(DA_1\) 所成的角为 \(90^\circ\)}
    {直线 \(BC_1\) 与 \(CA_1\) 所成的角为 \(90^\circ\)}
    {直线 \(BC_1\) 与平面 \(BB_1D_1D\) 所成的角为 \(45^\circ\)}
    {直线 \(BC_1\) 与平面 \(ABCD\) 所成的角为 \(45^\circ\)}
\end{problem}

\begin{problem}
已知函数 \(f(x)=x^3-x+1\)，则（\quad）

\choices
    {\(f(x)\) 有两个极值点}
    {\(f(x)\) 有三个零点}
    {点 \((0,1)\) 是曲线 \(y=f(x)\) 的对称中心}
    {直线 \(y=2x\) 是曲线 \(y=f(x)\) 的切线}
\end{problem}

\begin{problem}
已知 \(O\) 为坐标原点，点 \(A(1,1)\) 在抛物线 \(C:\ x^2=2py\)（\(p>0\)） 上，过点 \(B(0,-1)\) 的直线交 \(C\) 于 \(P\)，\(Q\) 两点，则（\quad）

\choices
    {\(C\) 的准线为 \(y=-1\)}
    {直线 \(AB\) 与 \(C\) 相切}
    {\(|OP|\cdot |OQ|>|OA|^2\)}
    {\(|BP|\cdot |BQ|>|BA|^2\)}
\end{problem}

\begin{problem}
已知函数 \(f(x)\) 及其导函数 \(f'(x)\) 的定义域均为 \(\R\)，记 \(g(x)=f'(x)\). 若 \(f\left(\frac32-2x\right)\)，\(g(2+x)\) 均为偶函数，则（\quad）

\choices
    {\(f(0)=0\)}
    {\(g\left(-\frac12\right)=0\)}
    {\(f(-1)=f(4)\)}
    {\(g(-1)=g(2)\)}
\end{problem}

\section{填空题}

\begin{problem}
\( \left(1-\frac{y}{x}\right)(x+y)^8 \) 的展开式中 \(x^2y^6\) 的系数为\fillinblank{}（用数字作答）.
\end{problem}

\begin{problem}
写出与圆 \(x^2+y^2=1\) 和 \((x-3)^2+(y-4)^2=16\) 都相切的一条直线的方程\fillinblank{}.
\end{problem}

\begin{problem}
若曲线 \(y=(x+a)e^x\) 有两条过坐标原点的切线，则 \(a\) 的取值范围是\fillinblank{}.
\end{problem}

\begin{problem}
已知椭圆 \(C:\ \frac{x^2}{a^2}+\frac{y^2}{b^2}=1\)（\(a>b>0\)），\(C\) 的上顶点为 \(A\)，两个焦点为 \(F_1\)，\(F_2\)，离心率为 \(\frac12\). 过 \(F_1\) 且垂直于 \(AF_2\) 的直线与 \(C\) 交于 \(D\)，\(E\) 两点，\(|DE|=6\)，则 \(\triangle ADE\) 的周长是\fillinblank{}.
\end{problem}

\section{解答题}

\begin{problem}
记 \(S_n\) 为数列 \(\{a_n\}\) 的前 \(n\) 项和，已知 \(a_1=1\)，数列 \(\left\{\frac{S_n}{a_n}\right\}\) 是公差为 \(\frac13\) 的等差数列.

\begin{enumerate}
\item 求 \(\{a_n\}\) 的通项公式；

\item 证明：\(\frac1{a_1}+\frac1{a_2}+\cdots+\frac1{a_n}<2\).
\end{enumerate}
\end{problem}

\begin{problem}
记 \(\triangle ABC\) 的内角 \(A\)，\(B\)，\(C\) 的对边分别为 \(a\)，\(b\)，\(c\)，已知 \(\frac{\cos A}{1+\sin A}=\frac{\sin 2B}{1+\cos 2B}\).

\begin{enumerate}
\item 若 \(C=\frac{2\pi}{3}\)，求 \(B\)；

\item 求 \(\frac{a^2+b^2}{c^2}\) 的最小值.
\end{enumerate}
\end{problem}

\begin{problem}
如图，直三棱柱 \(ABC-A_1B_1C_1\) 的体积为 \(4\)，\(\triangle A_1BC\) 的面积为 \(2\sqrt2\).

\begin{enumerate}
\item 求 \(A\) 到平面 \(A_1BC\) 的距离；

\item 设 \(D\) 为 \(A_1C\) 的中点，\(AA_1=AB\)，平面 \(A_1BC\perp\) 平面 \(ABB_1A_1\)，求二面角 \(A-BD-C\) 的正弦值.
\end{enumerate}

\bitmapfigure[max width=0.28\textwidth]{img/2022/new_gaokao_paper_1/q19_fig1.png}
\end{problem}

\begin{problem}
一医疗团队为研究某地的一种地方性疾病与当地居民的卫生习惯（卫生习惯分为良好和不够良好两类）的关系，在已患该疾病的病例中随机调查了 \(100\) 例（称为病例组），同时在未患该疾病的人群中随机调查了 \(100\) 人（称为对照组），得到如下数据：

\begin{center}
\begin{tabular}{|c|c|c|}
\hline
 & 不够良好 & 良好 \\
\hline
病例组 & 40 & 60 \\
\hline
对照组 & 10 & 90 \\
\hline
\end{tabular}
\end{center}

\begin{enumerate}
\item 能否有 \(99\%\) 的把握认为患该疾病群体与未患该疾病群体的卫生习惯有差异？

\item 从该地的人群中任选一人，\(A\) 表示事件“选到的人卫生习惯不够良好”，\(B\) 表示事件“选到的人患有该疾病”. \(\frac{P(B\mid A)}{P(\overline B\mid A)}\) 与 \(\frac{P(B\mid \overline A)}{P(\overline B\mid \overline A)}\) 的比值是卫生习惯不够良好对患该疾病风险程度的一项度量指标，记该指标为 \(R\).

\begin{enumerate}
\item 证明：\(R=\frac{P(B\mid A)}{P(\overline B\mid A)}\div \frac{P(B\mid \overline A)}{P(\overline B\mid \overline A)}=\frac{P(A\mid B)}{P(\overline A\mid B)}\cdot \frac{P(\overline A\mid \overline B)}{P(A\mid \overline B)};\)

\item 利用该调查数据，给出 \(P(A\mid B)\)，\(P(A\mid \overline B)\) 的估计值，并利用（1）的结果给出 \(R\) 的估计值.
\end{enumerate}

附：\(K^2=\frac{n(ad-bc)^2}{(a+b)(c+d)(a+c)(b+d)}\)，

\begin{center}
\begin{tabular}{|c|c|c|c|}
\hline
\(P(K^2\ge k)\) & 0.050 & 0.010 & 0.001 \\
\hline
\(k\) & 3.841 & 6.635 & 10.828 \\
\hline
\end{tabular}
\end{center}
\end{enumerate}
\end{problem}

\begin{problem}
已知点 \(A(2,1)\) 在双曲线 \(C:\ \frac{x^2}{a^2}-\frac{y^2}{a^2-1}=1\)（\(a>1\)） 上，直线 \(l\) 交 \(C\) 于 \(P\)，\(Q\) 两点，直线 \(AP\)，\(AQ\) 的斜率之和为 \(0\).

\begin{enumerate}
\item 求 \(l\) 的斜率；

\item 若 \(\tan \angle PAQ=2\sqrt2\)，求 \(\triangle PAQ\) 的面积.
\end{enumerate}
\end{problem}

\begin{problem}
已知函数 \(f(x)=e^x-ax\)，\(g(x)=ax-\ln x\) 有相同的最小值.

\begin{enumerate}
\item 求 \(a\)；

\item 证明：存在直线 \(y=b\)，其与两条曲线 \(y=f(x)\) 和 \(y=g(x)\) 共有三个不同的交点，并且从左到右的三个交点的横坐标成等差数列.
\end{enumerate}
\end{problem}

